\section{A quartic in forty-two unrestricted integers}
\label{sec:q42}
The new equation has \textbf{42 unrestricted integer unknowns}, exact degree
four, and 27 quadratic residuals. It improves the incoming 45-integer,
28-residual quartic without relying on the Jones universal expression or
Sun's ten-integer reduction. Its semantic source remains the inherited
native compiler and ZFC-refutation calibration. The new argument is a
cofinal scale restriction that orients a shared sign guard before any
power-of-two recovery.

\subsection{The complete quadratic system}
Use the 32 arithmetic variables
\[
\begin{gathered}
 J,F,\alpha,z,f,h,r,v,w,\tau,\eta,\zeta,y,\gamma,o,j,\\
 N,E,P,c,a,X,D,C_2,L,F_2,V,g,I,B,\Delta,u,
\end{gathered}
\]
one additional variable $q_h$, and three triples
$\mathbf t_A,\mathbf t_C,\mathbf t_K$. All42 range over $\mathbb Z$.
The letters $B,C_2,F_2$ are independent lifted registers, not radicals.
The four fixed numerals $(\beta,K_0,M_C,M_F)$ are the entire retained
package34 tuple, $\beta=2^{1064921}-1$; $M_F$ is already shifted.
Define only polynomial abbreviations
\[
 q=\beta J+1,\qquad k=\eta+\zeta,\qquad
 \operatorname{Pos}(x,y,t)=1+x^2+y^2+t(t+1),
\]
\[
 \widehat R=F(q+L-N)+J(M_C+qM_F)+q(N-L-q)+L\alpha.
\]
On the first two equations, $L=q^2-1$, $N=q^3$ and
$\widehat R=((q-1)(q-F)+\alpha)(q^2-1)+(M_C+qM_F)J$.
The 27 residuals are the following explicitly quadratic expressions:
{\small
\begin{align*}
e_0&=L-q^2+1,& e_1&=N-q(L+1),\\
e_2&=X-qw,& e_3&=a-E-(2v+1)N,\\
e_4&=E(X+1)-Xa,&e_5&=c-k(a-E)-\eta,\\
e_6&=P-E(c-\eta)-k,&e_7&=\tau^2-P^2+Pk-1,\\
e_8&=D-X-ac-(4a+3)\gamma,&e_9&=\Delta-a^2-4a-3,\\
e_{10}&=C_2-c^2,&e_{11}&=D^2-\Delta C_2-1,\\
e_{12}&=k-1-Eh-\widehat R,\\
e_{13}&=(K_0+w)(q-F-\alpha)-F-z(q-1),\\
e_{14}&=F_2-f^2,&e_{15}&=V+c-fo,\\
e_{16}&=V-cj+k-1-Eh,&e_{17}&=V^2+g-y^2,\\
e_{18}&=I-(2r+1)C_2,&e_{19}&=B-I^2,\\
e_{20}&=B-\Delta F_2+\Delta,&e_{21}&=y^2-Bg-1,\\
e_{22}&=w-4u^2,&e_{23}&=q-2q_h^2,\\
e_{24}&=F\alpha-\operatorname{Pos}(\mathbf t_A),\\
e_{25}&=F(q-F-\alpha)-\operatorname{Pos}(\mathbf t_C),\\
e_{26}&=\eta\zeta-\operatorname{Pos}(\mathbf t_K).
\end{align*}}
The single recommended equation is
\[
 \boxed{\quad Q_{42}=\sum_{j=0}^{26}e_j^2=0.\quad}
\]
The supplied exact expansion has391 symbolic terms. The coefficient of
$\tau^4$ is1, so specializing the fixed numerals cannot reduce its degree.
All variables occur; no positivity side condition is left implicit.

\subsection{What changes from the forty-five-variable source}
The predecessor guards $\eta$ and $\zeta$ separately using six ternary
positivity variables. We replace those equations by $e_{26}=0$, retaining
only three such variables. This initially establishes only that $\eta$ and
$\zeta$ have the same nonzero sign. At the same time, we replace its
$w=u^2$ by $w=4u^2$, adding no variable or residual. Hence the counts are
$45-6+3=42$ and $28-2+1=27$.

Every positive integer is represented by $\operatorname{Pos}$, by the
three-square theorem applied to $4n-3$. The restriction $w=4u^2$ is not
asserted for every predecessor solution: completeness will use the
canonical cofinal family. Omitting this strengthened scale leaves a real
negative-common-sign branch and does not give our proof.

\subsection{Recover the common sign before invoking the source theorem}
At a zero of $Q_{42}$ every $e_j$ vanishes. The dyadic guard gives $J>0$ and
$q\ge16$. The two simplex product guards, together with
$F+\alpha+(q-F-\alpha)=q>0$, give $F,\alpha,q-F-\alpha>0$.
The unchanged mask bounds give
\[
 q(q^2-1)<R:=\widehat R<q^4-2.
\]
They require no sign of $z$, $h$, $v$ or $k$.
The arithmetic lift implies
\[
 s=2v+1,\quad Y=sq^3=a-E,\quad E=XY,\quad
 c=kY+\eta,\quad a=(X+1)Y.
\]
Thus $Y\ne0$, $Z=|Y|\ge q^3$. If $u=0$, then $X=0$ and reducing the
main Pell norm modulo $H=4a+3$ gives $0\equiv1\pmod H$, impossible since
$|H|>1$. Therefore $u\ne0$, and the strengthened scale proves immediately
\[
 X\ge4q,\qquad |E|\ge4q^4>3R+1.                 \tag{Q1}
\]
This precedes the rank argument, so it is not inferred from an eventual
identity $X=2^R$.

Write $\sigma=\operatorname{sign}(k)$, $k_0=|k|$, $c_0=|c|$,
$A_0=|a+2|$ and $P_0=2XY^2+1$. Since $\eta\zeta>0$,
$0<|\eta|<k_0$; hence
\[
 k_0(Z-1)<c_0<k_0(Z+1),\quad c_0>k_0,\quad
 P_0>A_0>|E|>R.
\]
The consecutive-product and main Pell classifications give $n,p\ge1$ with
\[
 k_0=2\psi_{P_0}(n),\quad c_0=\psi_{A_0}(p),\quad p>n.
\]
The index equation and $P_0\equiv1\pmod{|E|}$ yield
\[
 2n\equiv\sigma(R+1)\pmod{|E|}.
\]
For $\sigma=1$, $n\ge(R+1)/2$. For $\sigma=-1$,
$n\ge(|E|-R-1)/2>R$ by(Q1). Both give $n\ge24$, so $p\ge25$ and
standard Pell growth proves
\[
 c_0>A_0\Delta^2,\qquad c_0>2p,\qquad c_0>2R.     \tag{Q2}
\]
The ordinary strong-rank lemma is therefore applicable in both sign
branches, not just after assuming $k>0$.

Put $f_0=|f|$ and $i=2r+1\ne0$. The retained strong equation gives a rank
$m$ with $f_0=\chi_{A_0}(m)$, $p\mid m$, $c_0\mid m$ and $m>2p$.
The auxiliary norm and the two exposed quotient equations give
\[
 V\equiv-c\pmod{f_0},\qquad V\equiv-R\pmod{c_0}.
\]
The odd-Chebyshev quotient identities and signed step-down give
$R\equiv\pm p\pmod{c_0}$; squaring the first congruence removes the sign
of $c$. The argument is unchanged for either sign of $Y$. The strict
bounds(Q2) force $p=R$. If $\sigma=-1$, this contradicts $p>n>R$.
Consequently $k>0$ and $\eta,\zeta>0$.

Only now invoke the predecessor's complete soundness theorem: keep its
arithmetic coordinates, put $u_{\rm old}=2u$, and represent the now-positive
$\eta$ and $\zeta$ by their old ternary guards. Every old residual vanishes
literally. The predecessor thus supplies a finite ZFC refutation. Its
additional negative-scale and signed-root analyses need not be repeated.

\subsection{Cofinal completeness and exact scope}
Conversely, its canonical refutation witnesses have $q=2^t$ with $t$ odd,
$R\equiv3\pmod4$, and $w=2^{R-t}$. Since $R>3q$, $R-t$ is even and at
least2. Choose $u=2^{(R-t-2)/2}$; then $w=4u^2$. Represent
$\eta\zeta>0$ by the new ternary guard and retain the other canonical
coordinates. All27 residuals vanish.

This proves both directions without deciding consistency. The retained
Q45, ordinary-rank, signed step-down and native-calibration sources are
identified and hashed in the archive. An independent symbolic expansion
checks the exact residual changes and degree. Those finite checks support
the literal expression; they do not replace the unbounded inherited
arithmetic/compiler theorems or constitute a new kernel certificate.

\subsection{Exact monic projections: thirty-six and thirty-two integers}
A residual $x-P(\mathbf y)=0$ with coefficient1 permits deleting $x$ and
substituting its integer polynomial value in every other residual. This
is an exact bijection of integer solution sets, with no divisibility
side condition. Acyclic chains of these substitutions preserve the same
consistency equivalence.

For the sextic, eliminate the six lifted coordinates
\[
 w=4u^2,\ C_2=c^2,\ F_2=f^2,\ B=I^2,\ L=q^2-1,
 \ E=a-(2v+1)N.
\]
The remaining21 residuals are cubic or lower. Their sum of squares is an
explicit polynomial in36 unrestricted integers of exact degree6;
the coefficient of $I^4g^2$ is1. The archive gives its435-term symbolic
expansion and the complete reconstruction map.

For the octic, eliminate instead
\[
\begin{gathered}
 w=4u^2,\ X=qw,\ C_2=c^2,\ \Delta=a^2+4a+3,\ F_2=f^2,\\
 B=I^2,\ g=y^2-V^2,\ L=q^2-1,\ N=q(L+1),\\
 P=E(c-\eta)+\eta+\zeta.
\end{gathered}
\]
The remaining17 residuals have degree at most4, so their sum of squares is
an octic in32 unrestricted integers. Its coefficient of $I^4V^4$ is1;
its exact expansion has395 symbolic terms. The substitutions $X=qw$ and
$N=q(L+1)$ are fully expanded after those for $w$ and $L$, respectively.
Neither eliminated register remains an existential unknown.

These are degree--variable tradeoffs, not42-variable quartics relabeled
with smaller counts. The supplied builders independently check the monic
identities, acyclic reconstruction, absence of eliminated symbols, degrees
and unit leading coefficients. Their semantic justification is the elementary
integer projection lemma above, composed with the proved Q42 reduction.
